\documentclass[11pt]{article}
\usepackage{amsmath, amssymb, amsthm}
\usepackage[margin=1in, top=1.2in, bottom=1.2in]{geometry}
\usepackage{hyperref}
\usepackage{parskip}
\hypersetup{colorlinks=true, linkcolor=blue, citecolor=blue, urlcolor=blue}
\theoremstyle{definition}
\newtheorem{definition}{Definition}[section]
\newtheorem{theorem}{Theorem}[section]
\newtheorem{lemma}{Lemma}[section]
\newtheorem{proposition}{Proposition}[section]
\newtheorem{corollary}{Corollary}[section]
\theoremstyle{remark}
\newtheorem{remark}{Remark}[section]
\newtheorem{example}{Example}[section]
\setlength{\parindent}{0pt}
\setlength{\parskip}{6pt}
\begin{document}

\begin{center}
{\Large\textbf{2.2 Measure Prerequisite-2}}\\
\vspace{0.3cm}
{\normalsize\today}
\end{center}
\vspace{0.5cm}


% --- Page 1 ---
Measurable function.

* If $(X, \mathcal{M})$ and $(Y, \mathcal{N})$ are measurable spaces a mapping $f: X \rightarrow Y$ is called $(\mathcal{M}, \mathcal{N})$-measurable or measurable if $f^{-1}(E) \in \mathcal{M}$ for all $E \in \mathcal{N}$.

Remark

Composition of measurable mappings are measurable.

Proposition

If $\mathcal{N}$ is generated by $\mathcal{E}$, then $f: X \rightarrow Y$ is $(\mathcal{M}, \mathcal{N})$-measurable iff $f^{-1}(E) \in \mathcal{M}$ for all $E \in \mathcal{E}$.

Proof: Only if condition is trivial. For if condition, observe $\mathcal{F} = \{E \subseteq Y | f^{-1}(E) \in \mathcal{M}\}$ is a $\sigma$-algebra, that contains $\mathcal{E}$ by hypothesis; it therefore contains $\mathcal{N}$.

Corollary

If $X$ and $Y$ are metric (or topological spaces),

% --- Page 2 ---
every continuous mapping is \((B_X, B_Y)\)
measurable.

* If \((X, \mathcal{M})\) is a measurable space, a real or
complex value \(f^n\) on \(X\) is called measurable if \(f\) is
\((\mathcal{M}, B_{\mathbb{R} \text{ or } \mathbb{C}})\) measurable.

\end{document}